\documentclass[10pt]{article}\usepackage[margin=0.8in]{geometry}\usepackage[T1]{fontenc}\usepackage{lmodern,amsmath,hyperref}
\title{Independent review: C0 Cantor reduction}\date{1 November 2026}\begin{document}\maketitle
Fresh isolated replay passed seven modules and 25 axiom readbacks: twelve main declarations, six producer controls and seven independent controls. The checked polynomial correspondence holds at its stated scope. The written divisor-class derivation is consistent under its smooth-curve assumptions and is separate from machine verification.
\section*{Native polynomial correspondence}
Native \texttt{three\_leads.rs:1015--1024} has norm $=-B$ and slope $=C$. Thus $C=0$ reaches the general two-Bezout fallback. With $U=u_1u_2$, $v_i=\ell+au_i$, $f-\ell^2=UQ$, and both Bezout identities, the literal numerator equals
\[\ell+U\{c_1a(e_1+e_2)+c_2(a^2+Q)\}.\]
The proof permits arbitrary $c_2$; an independent control uses $c_2=1$. The initial remainder is $\ell$ by its degree bound. The exact monic quotient is $X-x$, and the output remainder is $-\ell(x)$. This describes one polynomial degree-$4$ to degree-$1$ reduction, conditional on native egcd/division specifications; it is not their execution proof.
\section*{Written geometric fidelity}
On the smooth projective model of a squarefree monic quintic in characteristic not two, the ideals $(u_i,y-v_i)$ have quotient lengths two. Coprime supports separate the inputs. Away from ramification, $y+\ell$ is a local unit on the selected branch, so the multiplicities of $y-\ell$ equal those of $f-\ell^2$. At ramification, squarefreeness gives a simple root; local orders are $\operatorname{ord}(y)=1$ and $\operatorname{ord}(x-\alpha)=2$, so $y-\ell$ has order one. Repeated nonramified roots and residual/input overlap are retained with multiplicity.

At infinity, $y$ has pole order five and $\ell$ has pole order at most four. Hence
\[\operatorname{div}(y-\ell)=D_1+D_2+R-5O,\qquad
\operatorname{div}(x_{\mathrm{coord}}-x)=R+R'-2O.\]
The second equality includes ramified $R=R'$ with multiplicity two. In the actual quotient $\mathrm{Div}^0/\mathrm{Prin}$ these imply the output class $[R'-O]$, where $R'=(x,-\ell(x))$. No desired class-sum rule is assumed. The smooth curve, valuations and Picard carrier remain written mathematical inputs, not Lean constructions; original-field descent remains explicit.
\section*{Controls and provenance}
Independent controls cover nonzero $c_2$, wrong quotient, ramification factorization/simple coefficient, ordinate sign and repeated residual support. They are algebraic examples, not formal local-valuation proofs. All readbacks use only the standard permitted axioms.

Published input commit: \texttt{ceecc0b38f9a3e14370781b2275afc6e59b27894}. An initial archive-suffix error produced empty modules; that invalid attempt is preserved and superseded by nonempty, byte-compared extraction and a 25-declaration audit. An ambiguous control identifier was also corrected with its failure retained. Review was source-first, not blind. Cached dependencies remain trusted. No scientific execution or new costs were produced.
\end{document}
